\section{Conclusion}
To address the critical challenge of reducing training data requirements while preserving both rendering fidelity and semantic accuracy across static and dynamic scenes, we introduce a unified 3DGS backbone that jointly models semantic and temporal representations, together with an active learning framework based on Fisher Information for Next-Best-View selection. By quantifying the informativeness of candidate viewpoints with respect to geometric and semantic Gaussian parameters and the deformation network, our method identifies the most informative viewpoints to guide the training process. Our experimental results on Replica and Neu3D validate the effectiveness of the approach. These findings highlight that view selection guided by explicit information measures can lead not only to faster convergence but also to more reliable modeling of dynamic and semantic scenes. By reducing the reliance on large amounts of training data, our method enables more efficient and scalable training of high-quality 3DGS models.\\
Our framework offers a foundation for adaptive data acquisition. In robotics and SLAM, principled NBV selection can help agents actively choose viewpoints to improve mapping accuracy and semantic awareness under limited sensing budgets. In AR/VR, efficient view selection accelerates scene capture and dynamic updates. These impacts highlight the potential of our approach to enable efficient, scalable, and intelligent scene understanding across diverse domains.

\paragraph{Acknowledgment}
We greatly acknowledge financial support by the grants NSF FRR
2220868, NSF IIS-RI 2212433, and ONR N00014-22-1-2677.

% \paragraph{Broader Impacts}
% Our proposed Fisher Information-based view selection framework aims to lower the data and computational expenses with training high-quality semantic and dynamic 3DGS. Through efficient learning with fewer annotated views, our method can democratize access to high-fidelity 3D scene understanding for resource-constrained research and applications. In robotics, this opens the door to embodied agents that actively decide where to look in order to improve their understanding of a scene, supporting tasks such as autonomous navigation, manipulation, and exploration in previously unseen environments. In AR/VR, selective view acquisition can accelerate the creation of immersive digital twins and reduce the latency of real-time scene updates, which is crucial for collaborative or interactive experiences. Similarly, in digital heritage and remote inspection, our method can prioritize capturing the most structurally or semantically informative perspectives, allowing practitioners to reconstruct high-fidelity models under tight time or bandwidth constraints.\\
% More broadly, the framework promotes a shift toward active and informed learning in 3D scene understanding, where model performance is not limited by the raw volume of data but instead guided by principled measures of information. This perspective may inspire new directions in open-world SLAM, low-power edge perception, and multi-modal scene capture pipelines that integrate RGB, depth, and LiDAR. Taken together, these impacts suggest that our approach can serve as a stepping stone toward intelligent systems that learn efficiently, interact adaptively with their environment, and scale responsibly across a wide range of application domains.


% \paragraph{Limitations}
% While our method achieves strong performance across multiple datasets, it relies on several simplifying assumptions. First, we adopt a diagonal approximation of the Fisher Information matrix to make the computation tractable, which may ignore cross-parameter dependencies in high-curvature regions. Second, our evaluation is limited to indoor scenes with known camera intrinsics and well-calibrated multi-view setups, which may not generalize directly to outdoor or unstructured settings. Additionally, the method assumes that the pretrained 2D semantic, LSeg~\cite{Lseg} provides reliable and unbiased supervision, which might introduce errors in the 3D domain due to the domain gaps. 
% Future work could explore integrating our view selection framework into downstream tasks such as active exploration and open-world SLAM, where real-time informativeness estimation is crucial. Additionally, extending the Fisher Information metric to multi-modal inputs, like depth or LiDAR, may further enhance scene understanding in complex, open-world scenarios.